% ============================================================
% Proposition 4.4: Noise Sensitivity
% Status: FULL RECONSTRUCTION
% ============================================================

% --- Definitions and Assumptions ---
\begin{definition}[Noisy Estimation]
The estimated schema coherence is $\hat{\sigma}_A = \sigma_A + \epsilon$,
where $\epsilon \sim \mathcal{N}(0, \nu^2)$ is additive Gaussian
estimation noise. Similarly, the estimated relative depth is
$\hat{\delta}_A^{\text{relative}} = \delta_A^{\text{relative}} + \epsilon'$,
with $\epsilon' \sim \mathcal{N}(0, \nu^2)$ independent of $\epsilon$.
\end{definition}

\begin{definition}[Phase Assignment Error]
A phase assignment error occurs when the noisy estimates place the
agent in the incorrect phase relative to the ground-truth state:
\begin{itemize}[nosep]
\item \textbf{Schema error:} $\hat{\sigma}_A$ is on the opposite side
  of $\sigma_{\text{critical}}$ from $\sigma_A$.
\item \textbf{Depth error:} $\hat{\delta}_A^{\text{relative}}$ is on
  the opposite side of $\delta^*$ from $\delta_A^{\text{relative}}$.
\end{itemize}
\end{definition}

\begin{assumption}[Independence of Noise Sources]
The estimation noise for $\sigma_A$ and $\delta_A^{\text{relative}}$
is independent. This holds when these quantities are estimated from
different measurement instruments (Stage~2 OOD/ID ratio vs. depth
frontier estimation).
\end{assumption}

\begin{assumption}[Equal Noise Variance]
Both estimators have the same noise variance $\nu^2$. The general case
with $\nu_1 \neq \nu_2$ follows the same form with separate terms.
\end{assumption}

% --- Proposition Statement ---
\begin{proposition}[Noise Sensitivity]
\label{prop:noise}
Under Gaussian estimation noise $\hat{\sigma}_A = \sigma_A + \epsilon$
with $\epsilon \sim \mathcal{N}(0, \nu^2)$, the phase assignment error
rate is:
\[
P(\text{error}) = \Phi\!\left(-\frac{|\sigma_A - \sigma_{\text{critical}}|}{\nu}\right)
+ \Phi\!\left(-\frac{|\delta_A^{\text{relative}} - \delta^*|}{\nu}\right),
\]
where $\Phi$ is the standard normal CDF. For $\nu = 0.05$ and typical
distances from thresholds ($|\sigma_A - \sigma_{\text{critical}}| > 0.15$),
the error rate is $< 0.3\%$.
\end{proposition}

% --- Proof ---
\begin{proof}
We compute the probability that estimation noise causes a phase
misclassification.

\paragraph{Step 1: Schema phase error probability.}
Let $d_\sigma = |\sigma_A - \sigma_{\text{critical}}|$ be the margin
between the true $\sigma_A$ and the bifurcation threshold. Without
loss of generality, assume $\sigma_A > \sigma_{\text{critical}}$ (the
agent is in the schema-coherent phase). The correct assignment is
Phase~3. An error occurs when $\hat{\sigma}_A \leq \sigma_{\text{critical}}$.

Since $\hat{\sigma}_A \sim \mathcal{N}(\sigma_A, \nu^2)$:
\begin{align*}
P(\text{schema error}) &= P(\hat{\sigma}_A \leq \sigma_{\text{critical}}) \\
&= P\!\left(\frac{\hat{\sigma}_A - \sigma_A}{\nu} \leq \frac{\sigma_{\text{critical}} - \sigma_A}{\nu}\right) \\
&= P\!\left(Z \leq -\frac{d_\sigma}{\nu}\right) \\
&= \Phi\!\left(-\frac{d_\sigma}{\nu}\right),
\end{align*}
where $Z \sim \mathcal{N}(0,1)$ is the standard normal variable. By
symmetry, the same expression holds when $\sigma_A < \sigma_{\text{critical}}$
(the agent is in Phase~1/2) — the error probability depends only on
$d_\sigma = |\sigma_A - \sigma_{\text{critical}}|$.

\paragraph{Step 2: Depth phase error probability.}
By identical reasoning, let $d_\delta = |\delta_A^{\text{relative}} - \delta^*|$.
Under $\hat{\delta}_A^{\text{relative}} \sim \mathcal{N}(\delta_A^{\text{relative}}, \nu^2)$:
\[
P(\text{depth error}) = \Phi\!\left(-\frac{d_\delta}{\nu}\right).
\]

\paragraph{Step 3: Total error probability.}
Since the two noise sources are independent (Assumption~1), the events
``schema error'' and ``depth error'' are independent. The total phase
assignment error rate is the sum of the individual error rates:
\[
P(\text{error}) = P(\text{schema error}) + P(\text{depth error})
- P(\text{schema error} \cap \text{depth error}).
\]
The intersection term is second-order:
\[
P(\text{schema error}) \cdot P(\text{depth error})
= \Phi(-d_\sigma/\nu) \cdot \Phi(-d_\delta/\nu).
\]
For typical margins $d_\sigma, d_\delta > 0.15$ and $\nu = 0.05$,
each $\Phi(-3) \approx 0.00135$, so the product is $\sim 1.8 \times 10^{-6}$,
negligible compared to the individual terms. Dropping the intersection
term gives the stated bound:
\[
P(\text{error}) \leq \Phi(-d_\sigma/\nu) + \Phi(-d_\delta/\nu).
\]
Equality holds in the limit of small error probabilities (where the
intersection term is asymptotically negligible).

\paragraph{Step 4: Numerical example.}
For $\nu = 0.05$ and $d = 0.15$ (the minimum typical distance from
thresholds), $d/\nu = 3$. The standard normal tail probability is:
\[
\Phi(-3) \approx 0.00135.
\]
The total error rate is:
\[
P(\text{error}) < 2 \times 0.00135 = 0.0027 < 0.3\%.
\]
For $d = 0.20$ ($d/\nu = 4$), $\Phi(-4) \approx 3.2 \times 10^{-5}$,
giving $P(\text{error}) < 6.4 \times 10^{-5}$.
$\square$
\end{proof}

% --- Boundary Cases ---
\begin{boundary-case}
\textbf{$d_\sigma = 0$ (on the threshold):}
The error probability is $\Phi(0) = 0.5$ — random guessing. The phase
assignment is meaningless when $\sigma_A$ is exactly at the bifurcation
point. In practice, hysteresis (Proposition~4.3) prevents oscillation
in this degenerate case.

\textbf{$\nu \to 0$ (perfect estimation):}
$\Phi(-d/\nu) \to 0$ as $\nu \to 0$ for any $d > 0$. Error-free phase
assignment is achieved in the limit of infinite data.

\textbf{$\nu \to \infty$ (useless estimation):}
The error rate approaches $2 \times 0.5 = 1$ (minus the vanishing
intersection term). The estimate provides no information about the
true phase.

\textbf{Unequal noise variances ($\nu_1 \neq \nu_2$):}
The general form is:
\[
P(\text{error}) = \Phi\!\left(-\frac{d_\sigma}{\nu_1}\right)
+ \Phi\!\left(-\frac{d_\delta}{\nu_2}\right),
\]
where $\nu_1$ and $\nu_2$ are the standard deviations of the
$\sigma_A$ and $\delta_A^{\text{relative}}$ estimators respectively.
\end{boundary-case}

% --- Circular Dependency Check ---
\begin{circularity-check}
\textbf{Dependencies on earlier results:} The proof uses
Proposition~4.1 ($\sigma_{\text{critical}}$) and
Proposition~4.2 ($\delta^*$) to define the thresholds.
It does not depend on any later results.

\textbf{Forward dependencies:} The noise sensitivity analysis
informs the reliability requirements for Theorem~4.2
(faculty-validity correspondence) but is not logically required
by it.

\textbf{Self-circularity:} The proof assumes Gaussian noise with
known variance $\nu^2$. In practice, the noise distribution may be
non-Gaussian (e.g., heavy-tailed due to rare OOD samples). The
Gaussian assumption gives a closed-form error expression; for
non-Gaussian noise, the same qualitative behavior holds with
$\Phi$ replaced by the appropriate CDF. The numerical example
($0.3\%$ at $\nu = 0.05$, $d = 0.15$) is specific to Gaussian
noise and may differ under heavy-tailed noise. The inequality
$P(\text{error}) \leq \text{sum of tail probabilities}$ holds
for any noise distribution by the union bound.
\end{circularity-check}
